% Physical unit 032; accepted Session 03 after fresh 400-dpi authority replay.
\GordanUnitStart{032}{printed page 27}
\GordanTopFolio{27}
\begingroup\small
ex formulis:
\[
\begin{array}{rcll}
 \sqrt{1-e^2}\,\tg B_0&=&\tg l_0, &\text{Cap. I \S\ 7}\vspace{0.2em}\\
 {2\over\pi}{\Sigma(l_0)\over b}&=&1+o\sin^2l_0-{3\over4}o^2\sin^4l_0+{5\over4}o^3\sin^6l_0\ldots, &\text{Cap. III \S\ 2}\vspace{0.2em}\\
 {a\over b}\left\{1-{2\over\pi}\mathcal A(l_0)\right\}&=&2p\cos l_0+\left\{1-2p+{25\over8}p^2\right\}p^2\cos^3l_0\ldots, &\text{Cap. III \S\ 4}\vspace{0.2em}\\
 x&=&e\sin B_0, &\text{Cap. II \S\ 9}\vspace{0.2em}\\
 x'&=&{\cos B_0\over \cos l_0}. &
\end{array}
\]

\sect{4}

Restat, ut computentur constantes:
\[
 q\qquad x\qquad E\qquad v.
\]
Quam ad rem formulam adhibere placet hanc:
\[
 \sqrt{x'}={1-2q+2q^4\cdots\over 1+2q+2q^4\cdots}\qquad \text{fund. pag. 184.}
\]
Quum $q$ quantitas minima sit, quod quidem ad terram attinet, nunquam numerum \gq{0,0004} excedens, $q^4$ omnino negligi licet:

Est igitur:
\[
 \sqrt{x'}={1-2q\over 1+2q}
\]
vel:
\[
 2q={1-\sqrt{x'}\over 1+\sqrt{x'}}
 ={\sqrt{\cos l_0}-\sqrt{\cos B_0}\over \sqrt{\cos l_0}+\sqrt{\cos B_0}},
\]
unde $q$ facile computari potest.

Qua in re, cum interdum $x'$ non accurate definitum sit, fieri posse concedimus, ut $q$ non extremis numeris computari possit.

Quod ne fiat, explicanda est alia formula.

Est:
\[
 \sqrt{x}=2\sqrt[4]{q}\,{1+q^2\over 1+2q}
 \qquad \text{fund. pag. 184}
\]
\[
 q={x^2\over16}\left\{{1+2q\over 1+q^2}\right\}^{4}
\]
\[
 \log q=\log{x^2\over16}-4\log(1+q^2)+4\log(1+2q)
 =\log{x^2\over16}+8q-12q^2+{8\over3}q^3.
\]
Sin $q$ accuratius computandum est, computetur oportet quantitas $q_1$ aliqua, propius ad $q$ accedens, e formula
\[
 \log q_1=\log{x^2\over16}+{x^2\over2}
\]
quo facto $q$ definitur ex aequatione
\[
 \log q=\log{x^2\over16}+8q_1-12q_1^2+{8\over3}q_1^3.
\]
\endgroup
